\section{Joint contraction and bounded positive realizations}
\label{sec:r74-joint-contraction}

The reversible classification has a complementary dissipative regime. The
condition below is imposed on all command products through one supplied norm,
and the construction gives one machine for all words. A separate realization
for each command would not establish that conclusion. Invariant-polytope
positive realization is classical \cite{Heller65,MonrasWinter14}; the point of
the statement is an explicit sufficient certificate with bounded probability
decoders, a prescribed pure start, and a stated number of live labels.

Let $A$ and $C$ be nonempty finite sets, $a=|A|$, and let
\begin{equation}\label{eq:r74-contracting-response}
 b(w,c)=p_c+\ell_c(x_w),\qquad
 x_{wa}=F_ax_w,\qquad x_\varnothing=x_0\in\R^n,
\end{equation}
where $F_a$ are linear, $p\in\Delta(C)$, $\sum_c\ell_c=0$, and the displayed
values are probabilities for every word. For a norm $\|\cdot\|$ write
$\|\ell\|_*$ for its dual norm. Suppose we are given $0<\rho<1$ and $r>0$
such that
\begin{equation}\label{eq:r74-contraction-certificate}
 \|F_ax\|\le\rho\|x\|\quad(a\in A),\qquad
 r\|\ell_c\|_*\le p_c\quad(c\in C).
\end{equation}
In particular, strictly positive $p_c$ give such an $r$ for every finite
collection of output functionals. Zero components are allowed only when the
corresponding functional vanishes. The second condition certifies that the
entire ball $rB$, not only the actual orbit, has legal decoder rows.

\begin{theorem}[A common contracting realization]\label{thm:r74-contraction}
Under \eqref{eq:r74-contraction-certificate}, the response
\eqref{eq:r74-contracting-response} has an exact stochastic realization from
a prescribed pure label, valid for every command word, with at most
\begin{equation}\label{eq:r74-contraction-count}
 \sum_{j=0}^{T-1}a^j+
 \left\lceil\left(1+\frac{2}{1-\rho}\right)^n\right\rceil
 \quad\hbox{live labels},
\end{equation}
where $T\ge1$ may be any integer satisfying
$\rho^{T-1}\|x_0\|\le r$. For $x_0\ne0$ one may take
\[
 T=1+\max\left\{0,
 \left\lceil\frac{\log(\|x_0\|/r)}{\log(1/\rho)}\right\rceil\right\};
\]
for $x_0=0$, the response is constant and one label suffices. No minimality or
computational-efficiency assertion is made for this upper bound.
\end{theorem}
\begin{proof}
Put $\varepsilon=1-\rho$. Choose a maximal $\varepsilon$-separated finite
set $V$ in the unit ball $B$, where distinct selected points have distance
greater than $\varepsilon$. The open balls of radius $\varepsilon/2$ about
these points are pairwise disjoint and lie in $(1+\varepsilon/2)B$.
Lebesgue volume, which scales by the $n$th power, gives
\[
 |V|\le (1+2/\varepsilon)^n.
\]
This bound also proves that the greedy selection stops; maximality says
that $V$ is an $\varepsilon$-net of $B$. If $P=\operatorname{conv}(V)$,
then $\rho B\subseteq P\subseteq B$. Indeed, for every real linear
functional $f$, choose a maximizer $x\in B$ and a point $v\in V$ at
distance at most $\varepsilon$ from $x$. Symmetry of the unit ball gives
\[
 \max_{v\in V} f(v)\ge f(x)-\varepsilon\|f\|_*
 =(1-\varepsilon)\|f\|_*.
\]
The separating-hyperplane characterization of convex containment proves
$\rho B\subseteq P$. Consequently
\[
 F_a(rP)\subseteq\rho rB\subseteq rP
 \quad(a\in A).
\]

Use one recurrent label for each $v\in V$, representing the vector $rv$.
For every such label and command choose nonnegative barycentric coefficients
of $F_a(rv)$ in $rV$; they sum to one and define a row of $Q_a$. Give that
label decoder $p_c+\ell_c(rv)$. The decoder is nonnegative by the second
condition in \eqref{eq:r74-contraction-certificate}, and its entries sum
to one because $p$ has mass one and $\sum_c\ell_c=0$.

In addition use one transient label for each word $w$ with $|w|<T$,
representing its actual vector $x_w$. Its decoder is the actual probability
row $b(w)$, legal by hypothesis. If $|w|<T-1$, command $a$ sends it
deterministically to the label $wa$. If $|w|=T-1$, then
\[
 \|F_ax_w\|\le\rho^T\|x_0\|\le\rho r,
\]
so $F_ax_w$ is in $rP$. Choose its barycentric row in the recurrent labels.
Start at the single empty-word label. There are exactly
$\sum_{j=0}^{T-1}a^j$ such transient labels before any duplicate states
are merged. Induction on the word shows that the probability-weighted
represented vector is always $x_w$: the transition identity holds on
each individual label, and linearity preserves it under a mixture. Pairing
with all $\ell_c$ and adding $p$ gives the entire response. This proves
the count, including the pure initial label and every first command.
\end{proof}

The assumption may also be supplied as a uniform product bound. If a known
norm $\|\cdot\|_0$ and constants $K\ge1$, $0<q<1$ satisfy
\[
 \|F_w\|_0\le Kq^{|w|}\qquad(w\in A^*),
\]
choose any $q<\rho<1$ and set
\[
 \|x\|'=\sup_{w\in A^*}\rho^{-|w|}\|F_wx\|_0.
\]
This is a norm with $\|x\|_0\le\|x\|'\le K\|x\|_0$ and
$\|F_ax\|'\le\rho\|x\|'$. Hence
Theorem~\ref{thm:r74-contraction} applies with
$\|x_0\|'\le K\|x_0\|_0$ and any $r$ satisfying
$r\|\ell_c\|_{0,*}\le p_c$. This is a consequence of a supplied uniform
certificate, not a decision procedure for joint spectral radius. A common
Euclidean contraction certificate, for example, is the finite collection
$F_a^{\mathsf T}HF_a\preceq\rho^2H$ for a supplied positive definite $H$.
The actual net and barycentric rows provide finite implementation data;
their exact algebraic or sampling encodings remain distinct resources.

